\section{Тест производительности}

Тест производительности сравнивает время работы реализованного алгоритма на основе суффиксного дерева с наивным подходом перебора всех циклических ротаций с использованием STL функций. Наивный метод генерирует все $n$ возможных ротаций и находит минимальную через \texttt{std::min\_element}.

Тестовые данные генерируются Python-скриптом с различными параметрами: короткие строки (до 1000 символов), средние (до 10000 символов) и длинные (до 100000 символов). Время измеряется в микросекундах (\texttt{us}) с использованием библиотеки \texttt{<chrono>}.

\subsection{Результаты базовых бенчмарков}

\begin{lstlisting}[language=bash]
artems@MacBook-Air-artem-3 lab5 % ./benchmark

String length: 1000
-----------------------------------
Suffix Tree: 831 us
Naive std::min_element: 952 us
-----------------------------------
Speedup: 1.14561x

artems@MacBook-Air-artem-3 lab5 % python3 test_generator.py 10000 random | ./benchmark

String length: 10000
-----------------------------------
Suffix Tree: 3486 us
Naive std::min_element: 25357 us
-----------------------------------
Speedup: 7.27395x

String length: 50000
-----------------------------------
Suffix Tree: 25256 us
Naive std::min_element: 379677 us
-----------------------------------
Speedup: 15.0331x
\end{lstlisting}

\subsection{Анализ результатов}

\subsubsection{Растущее преимущество с увеличением размера}

Результаты демонстрируют, что разница в производительности растет экспоненциально с увеличением длины строки:

\begin{enumerate}
    \item \textbf{Строки длины 100:} Suffix tree примерно такой же по производительности, что и наивный подход (1.14x быстрее)
    \item \textbf{Строки длины 1000:} Преимущество возрастает до 7x 
    \item \textbf{Строки длины 50000:} Превосходство достигает 15x
\end{enumerate}

Это объясняется асимптотической сложностью:
\begin{itemize}
    \item \textbf{Suffix Tree:} $O(n)$ построение + $O(n)$ поиск + $O(1)$ вхождение = $O(n)$
    \item \textbf{Naive approach:} $O(n)$ ротаций $\times$ $O(n)$ сравнение $= O(n^2)$
\end{itemize}

\subsubsection{Влияние повторяющихся паттернов}

Специальный тест на строке, состоящей из повторяющегося паттерна:

\begin{lstlisting}[language=bash]
Test 4: Repetitive pattern (aaabaaabaaab...)
artems@MacBook-Air-artem-3 lab5 % python3 test_generator.py 50000 worst_case | ./benchmark

String length: 50000
-----------------------------------
Suffix Tree: 12060 us
Naive std::min_element: 477186 us
-----------------------------------
Speedup: 39.5677x
\end{lstlisting}

Даже на повторяющихся данных suffix tree сохраняет линейную сложность, в то время как наивный подход остаётся $O(n^2)$.

\subsubsection{Критическая точка пересечения}

Анализ показывает, что реализация suffix tree ощущается намного выгоднее наивного подхода начиная с длины строки примерно 200 символов. Для строк меньшего размера наивный подход может быть достаточным, учитывая простоту реализации.

\pagebreak
